\textbf{Transparencia de distribución rota.} Para quien usa S3, un GET o un PUT es una sola llamada. Por dentro intervienen al menos dos subsistemas: el de índice, que guarda los metadatos y la ubicación de todos los objetos de la región, y el de ubicación (\textit{placement}), que asigna almacenamiento a los objetos nuevos y no funciona sin el índice~\cite{aws-s3-2017}. La transparencia de fallos pide ocultar al cliente el fallo y la recuperación de un recurso~\cite{vansteen2023}, y eso fue lo que no se cumplió el 28 de febrero de 2017. Los dos subsistemas tuvieron que reiniciarse por completo y, mientras tanto, S3 no atendió peticiones en US-EAST-1. El corte se vio también en servicios que guardan datos en S3: la consola de S3, el lanzamiento de instancias nuevas en EC2, los volúmenes EBS que necesitaban datos de una instantánea y Lambda~\cite{aws-s3-2017}. Hubo una dependencia que desde fuera nadie veía. La consola con la que AWS administra su panel de estado usaba S3, así que hasta las 11:37 PST no pudo actualizar el estado de cada servicio y avisó por Twitter~\cite{aws-s3-2017}.

\textbf{Teorema CAP.} Gilbert y Lynch demuestran que un servicio no puede garantizar consistencia y disponibilidad a la vez cuando la red puede perder mensajes entre nodos~\cite{gilbert2002}. En este incidente no hubo partición. El informe no habla de pérdida de comunicación entre nodos; lo que faltó fue capacidad, porque se retiraron servidores de más~\cite{aws-s3-2017}. Por eso no lo presentamos como un ejemplo de elegir entre C y A. Sí vemos una decisión parecida en la recuperación: antes de volver a atender, S3 ejecutó las comprobaciones que validan la integridad de los metadatos, y el informe reconoce que eso tardó más de lo esperado~\cite{aws-s3-2017}. Nuestra lectura es que se aceptó más tiempo caído a cambio de no servir un índice sin verificar.

\textbf{Tolerancia a fallos.} Avizienis et al. separan la falta, el error y el fallo de servicio~\cite{avizienis2004}. Aquí la falta fue humana y ocurrió en operación: un integrante autorizado del equipo escribió mal uno de los parámetros de un comando. El error fue el estado en que quedó el sistema, con el índice y el \textit{placement} sin buena parte de su capacidad. El fallo llegó cuando la API dejó de responder. Para el cliente es un fallo de omisión, porque el servidor no contesta las peticiones que le llegan~\cite{vansteen2023}. AWS afirma que S3 está diseñado para perder capacidad significativa con poco o ningún impacto~\cite{aws-s3-2017}. Esta vez no alcanzó, por dos motivos que da el mismo informe: la herramienta dejó retirar demasiada capacidad demasiado rápido, y esos subsistemas llevaban muchos años sin reiniciarse por completo en las regiones grandes.
